\subsection{Generation by the individual disks}\label{sec:diskgeneration}

The morphism calculation leaves open whether the disks generate. A Weinstein handle of middle index has a distinguished Lagrangian disk, its cocore, obtained by flowing the local unstable disk outwards. The handle-attachment theorem gives generators in terms of these cocores. We shall choose a Liouville primitive whose unstable manifolds are precisely the geometric disks used in Proposition~\ref{prop:diskhoms}. We work in the idempotent-complete derived partially wrapped category.

\begin{theorem}\label{thm:diskgeneration}
For the fixed Morse presentation of Proposition~\ref{prop:diskhoms}, the individually labelled common-path thimbles form a full exceptional collection, in increasing terminal-angle order, in the category stopped at the core of the marked regular fibre. The group permutes their isomorphism classes, and within each common critical value any ordering is allowed. The subcategory generated by the disks over a value with $l$ critical points is equivalent to $\mathrm{Perf}(\C)^l$.
\end{theorem}

\begin{proof}
Choose a separation by Lemma~\ref{lem:samedisks} and use the flat completion and cylindrical representatives in the preceding proof. We first prepare a regular region over which the fibration is a product, including a corridor to the marked ray. The connection on this region can have curvature initially.

\medskip
\noindent\emph{A product away from the vanishing paths.}
Enclose each path in a thin closed region formed by a disk about its starting point and a strip along the path widening to a small angular sector on its radial end. Take the regions disjoint from one another and from the marked direction at infinity. In a large base disk $B_R$ the closed middle $M$ left after deleting their interiors is a disk with corners disjoint from every path. Enlarge it slightly across its edges and beyond the outer circle, and smooth to obtain an open disk $V$ avoiding all the paths. Contract $V$ to $q_0=R$ on the marked ray by a smooth contraction $h_t$ fixing $q_0$. For a cutoff $\chi$ supported in $V$ and equal to one on a neighbourhood of $M$, put $k_s(w)=h_{s\chi(w)}(w)$, extended by the identity. Thus $k_s$ changes points only through regular values, fixes the paths with their neighbourhoods, and $k_1=q_0$ near $M$.

Let $\Lambda$ be the flat-end primitive, adapted to the individual disks, and $\Theta=\Lambda-\pi^*\alpha_0$, using the base coordinate in which the plane primitive is standard. On the product fibre end and in the product chart at $q_0$ we have $\Theta=\lambda_F$. Choose a nonnegative plane area form $\omega_b$ allowing zeros, zero on disks about the critical values and standard outside slightly larger disks disjoint from $\overline V$. Choose a primitive $\beta$ vanishing on the smaller disks and equal to $\alpha_0+c\,d\theta$ near infinity. Indeed the constant $2\pi c$ is the integral of $\omega_b-d\alpha_0$. After removing this period the remaining difference on the end is exact, as are the restrictions on the inner disks. Cutoffs of these exact corrections give the required primitive. At large $R$ its angular coefficient is positive.

On the pullback
\[
 E_s=\{(w,x):w\in B_R,\ \pi(x)=k_s(w)\},\qquad \pi_s(w,x)=w,
\]
take the primitive $\operatorname{pr}_E^*\Theta+\pi_s^*(\alpha_0+A\beta)$ and truncate in the product fibre end. The pullback is smooth even at $s=1$: on the changed region $\pi$ is a submersion, while near each critical point $k_s$ is the identity. Horizontal lifting along $s\mapsto k_s(w)$ identifies the truncated family, by the identity near the paths and in product coordinates near the page. A small neighbourhood of $q_0$ stays in its product chart for all $s$ by compactness.

The new horizontal lift of $v$ is $(v,\widetilde{dk_s(v)})$. Its horizontal area form is $d\alpha_0+A\omega_b+k_s^*((d\Theta)_{\rm hor})$. On the changed compact part the last term is bounded independently of $A$ and on the fibre end it vanishes. Thus a single large $A$ gives positive area even when $dk_s$ drops rank. Elsewhere one has added nonnegative base area, leaving the standard critical charts intact. On the horizontal face the fibre field is outward. On the vertical face the primitive on the angular lift has the positive summand $A\beta(\partial_\theta)$ and the remaining terms are uniformly bounded. At $s=0$ first turn on $A\beta$; its angular contribution at $R$ is nonnegative throughout. The disks stay Lagrangian and conical, since the added base primitive vanishes on the radial tails. The far radial momentum is now $p+A(p+c)$ for $p=|w|^2/2$. This increasing change of radial coordinate preserves the page direction and fibre core.

The interpolation in $k_s$ keeps the disk neighbourhoods, the page germ and the horizontal product collar fixed under these identifications. Relative contact isotopy on the remainder of the compact boundary, followed by compact Moser transport, identifies the completed pairs and their disks. Here the contact conformal factor is absorbed by translation in Liouville time on the collars before applying Moser, and the transports may be fixed on the unchanged neighbourhoods. At the endpoint freeze radial transport on a short outer annulus as in Proposition~\ref{prop:diskhoms}, resampling the term from $\Theta$ and retaining $\alpha_0+A\beta$. The horizontal curvature and its resampled error are uniformly bounded there independently of $A$, since base addition preserves the connection. Choose $A$ at the outset to dominate this bound as well. The outer middle and its edge collars are already products, since $k_1$ is constant there on an open neighbourhood. In that part freezing preserves $\lambda_F$ and the product extends over the infinite middle. On each disk the radial substitution and collar identification again preserve the cylindrical representative.

\medskip
\noindent\emph{A Weinstein primitive adapted to the disks.}
Straighten the separated paths to radial rays $r e^{i\theta_i}$, $r\ge1$, moving their caps into $r>3/4$ and keeping their angles at infinity. This uses a compactly supported orientation-preserving base diffeomorphism. Prescribe translations and rotations on path tubes near their starts and the identity near an outer circle. The complementary regions, slit along those tubes, are disks. Their boundary identifications extend over collars and then across the disks. Such a postcomposition preserves fibres and horizontal transport on corresponding paths. Rotations at the centres are absorbed by unitary scalar changes in the quadratic charts. The disk $r\le2/3$ and a corridor about the entire marked ray now have the same product fibre. The projection, again denoted by $\pi$, is $c_i+\sum_jz_j^2$ in a standard chart, with $c_i=e^{i\theta_i}$.

Subtract a base one-form equal to the base part of the total primitive on a slightly smaller product disk and corridor, equal to its base part at infinity, and zero near the centres. These prescriptions agree on overlaps and extend by cutoffs. Write $\vartheta$ for the difference. On a smaller Morse ball a local exact change gives $\vartheta=\lambda_{\rm std}=\sum_j(x_jdy_j-y_jdx_j)/2$. Choose a generic Weinstein fibre primitive $\lambda_F^W=\lambda_F+dg$ with $g$ compactly supported. Indeed the strictly plurisubharmonic fibre exhaustion has no critical points beyond a compact level. Perturb it to Morse on the interior, keeping the regular collar, and pull the new forms back by compactly supported exact Moser transport. The resulting primitive differs by the differential of a compactly supported function. A finite fibre needs no perturbation. Add $d(\chi g)$ using a base cutoff equal to one on $r\le0.55$ and on a smaller marked corridor and supported inside the product. The complete page ribbon follows the change of its restricted primitive, away from the links of the disks at infinity \cite[Example~3.11]{GPS}.

The restriction of $\vartheta$ to each disk is exact since base area restricts to zero and the disk is contractible. It already vanishes near the centre, where the disk is a rotated real plane. Normalize its primitive to zero there and subtract the differential of a tubular extension cut off near the disk. Use disjoint tubes away from the smaller product disk and marked corridor. On a radial tail the fibre slice is compact and the gauge can be chosen independent of $p$, with compact fibre support. This allows both restriction constants on a zero-dimensional slice. We have obtained
\[
  \vartheta|_{TD_i}=0,\qquad \vartheta=\lambda_{\rm std}
                     \quad\text{near each critical point}.
\]
On the smaller product disk and the marked corridor, $\vartheta=\lambda_F^W$. Use $\Lambda_K=\vartheta+K\pi^*\alpha_0$ for large $K$. To deform to it, first add the large base term, then interpolate the subtractions. On regular compact sets the curvature errors and angular covectors are bounded independently of $K$. In the Morse charts the subtracted base form was zero. The connection and fibre forms are preserved. At large base radius the angular gauges have bounded derivatives with compact fibre support, giving strictly positive radial growth. The product fibre end still has $\dot I=I$ in its symplectization coordinate. These estimates give complete convex enlargements. On the far disk collars the initial and final adapted primitives vanish along disk tangents, as do all base forms along radial tangents. The disks thus persist as exact cylindrical objects during this deformation.

\medskip
\noindent\emph{The ascending manifolds.}
Set $H=|\pi|^2$ and $\iota_{Z_K}d\Lambda_K=\Lambda_K$. On compact regular subsets with $r$ bounded away from zero the projected field tends to $w/2$ by the horizontal two-plane formula. Hence $dH(Z_K)>0$ there for large $K$. The product fibre end gives the same bound for bounded base radii away from zero. At base infinity $\dot p=p+O(1/K)$ in the flat angular charts. Close to a critical point write $w=c+q(z)$ with $q=\sum_j z_j^2$, and put $a=4|z|^2$. On the standard horizontal complex line $d\pi$ is conformal of scale $\sqrt a$. The standard Liouville field is half the radial field on the $z$-coordinates, so its projection is $q$. The new projection is the average of $q$ and $w/2$ with weights $1$ and $Ka$. Thus
\[
                  dH(Z_K)=
           \frac{2\operatorname{Re}(\bar w q)+Ka|w|^2}{1+Ka}.
\]
Here $|q|\le|z|^2$ and $|w|$ is bounded away from zero. For fixed large $K$ the right side is bounded below by a positive multiple of $|z|^2$ in a small ball off its centre. At the centre $Z_K=0$ and $H$ has index $n=\dim_\C X$ by its quadratic part $2\operatorname{Re}(\bar c q)$. Rotate to $c=1$ in this calculation. The added linear covector is $K\,d\operatorname{Im}q/2$, so the linearized field has eigenvalues $1/2+K$ on the real plane and $1/2-K$ on the imaginary plane. The former is tangent to the outgoing disk.

To get a proper Weinstein function, choose a nonnegative radial function $h_\infty(I)$ on the product fibre end, zero on a large compact fibre part, nondecreasing and eventually linear and unbounded. Choose the fibre Lyapunov function $h_0\ge0$ for $\lambda_F^W$ to have critical values below $1/100$, to agree with $h_\infty$ far out, and to be uniformly close to it. Indeed one can add $\varepsilon$ times a shifted Morse Lyapunov function to $h_\infty$ on the compact region, extend increasingly along the collar until $h_\infty'>0$, and cut off the summand where this derivative is bounded below. With a radial base cutoff $\rho$ equal to one near zero and zero for $r\ge1/3$, put
\[
                         \Phi=H+h_\infty+\rho(r)(h_0-h_\infty).
\]
This is proper and near zero is a product Weinstein function. On the transition its fibre derivative is a convex combination of nonnegative derivatives, and the cutoff error is dominated by $dH(Z_K)=r^2$ there for small $\varepsilon$. Elsewhere $dh_\infty(Z_K)\ge0$ by the product fibre-end equation. At the centres the preceding quadratic bound gives the gradient-like inequality since both $dH$ and $Z_K$ are $O(|z|)$. The critical points besides the centres are the fibre critical points over zero, all of index at most $n-1$. The centres all have value one, so no trajectory connects two distinct centres.

The sublevel $\Phi\le1/5$ lies inside $r<0.55$ with primitive $\lambda_F^W+K\alpha_0$. Its completion is the entire stabilization $F\times\C$: the same product exhaustion extends throughout that product, with no critical values above $1/5$, so backward flow from outside reaches the level. Truncate $F$ above its core and below the start of $h_\infty$. The product of this truncation with $w=\sqrt{1/5}$ is the incoming page ribbon. On the whole marked corridor $Z_K=(Z_F^W,w/2)$, and forward flow preserves the reference direction and identifies the incoming with the original stopped core. Every point in that corridor with nonzero base radius flows outward at its fixed angle. The stable attaching spheres of the critical points miss the corridor.

The elementary Weinstein cobordism above this sublevel is given by disjoint critical handles on the free part of the incoming boundary \cite[\S6.3, Lemma~6.6]{GirPar}. Each prescribed $D_i$ is its completed unstable disk. Since the path is radial, $\Lambda_K|_{TD_i}=0$ and $Z_K$ is tangent everywhere. The real disk at the critical point has the expanding tangent plane, hence is the local unstable manifold by uniqueness. Every point of $D_i$ flows backwards into that germ, since each bounded segment is compact and $dH(Z_K)$ has a positive minimum on closed segments away from the centre. Conversely its forward saturation lies in the disk. One may keep this Lagrangian germ in the relative handle normal form: in the chart the Darboux interpolation for the added area uses the primitive $K(q_R\,dq_I-q_I\,dq_R)/2$, $q=q_R+iq_I$, which vanishes on the real plane. The Moser field is tangent there and has zero linearization. In the resulting coordinates the given and standard primitives restrict to zero on the plane, so a primitive of their exact difference can be set to zero there. The cutoff in the relative normal form preserves this restriction. Flow continues the identified cocore through the regular levels with the disk's transported grading.

\medskip
\noindent\emph{The retained page.}
The incoming stopped product is represented by the sector $P=F\times\C_{\operatorname{Re}\ge0}$ in \cite[(3.3)]{GPS}. A sufficiently small ribbon for the page core misses all attaching spheres by the corridor construction. Write $\mathcal L$ for their union and $P'$ for the attached sector. Under the attachment functor $\mathcal W(P,\mathcal L)\to\mathcal W(P')$ a small linking disk at a component of an attaching sphere maps to its handle cocore up to a grading shift. Both disks are images of a cotangent fibre of the handle overlap \cite[\S11.3]{GPS}. The disjoint charts allow all the identifications simultaneously. In dimension one an attaching sphere is $S^0$. The two overlap fibres on the handle side are cotangent fibres over the same interval and both give the one cocore up to shift. The source is generated by the linking disks of $\mathcal L$ by \cite[Theorem~1.16]{GPS}.

The original fibre page is still a stop after the attachment. It therefore contributes its own linking disks to the generation theorem, and we must see where these disks come from. For this purpose use the Weinstein convexification of the target with its original page stop. The incoming stabilization has handles only of index less than $n$. The generic fibre core is a mostly Legendrian stop, and the new cocores avoid it at infinity. By \cite[Theorem~1.13]{GPS} the target is generated by the new cocores and linking disks to the smooth Legendrian part of the original page core. These latter disks are already in the image of the attachment functor. Choose each in the incoming product corridor near its page and away from $\mathcal L$. The sector-to-ribbon and ribbon-to-core comparisons with the disjoint extra stop give an object of $\mathcal W(P,\mathcal L)$. Its page neighbourhood survives attachment, so the image is the corresponding outgoing linking disk, as in \cite[(9.6)]{GPS}. Thus the images of the attaching linking disks generate the target in the perfect envelope. They are exactly the cocores identified with our thimbles. Proposition~\ref{prop:diskhoms} supplies their morphism groups and the action on the initial collection. The endomorphism algebra of the sum of $l$ mutually orthogonal exceptional objects has cohomology $\C^l$ concentrated in degree zero, giving $\mathrm{Perf}(\C)^l$ as asserted.
\end{proof}
